\section{Threat Model}
\label{sec:threat-model}

\paragraph{What the attacker controls.}
The attacker controls one agent in the account. This is the realistic starting
point for an agent platform: an agent processes input it did not write, so an
indirect prompt injection, a poisoned tool result, or a dependency in the
agent's own image is enough to run attacker-chosen code with that agent's
credentials. We do not model how the first agent is taken; we start from the
point where attacker code runs as agent~A.

\paragraph{What the attacker wants.}
Something agent~A is not authorised to reach, held by agent~B in the same
account. In the terms of Equation~\ref{eq:esc}, the attacker wants an element
of $\eff(A) \setminus \auth(A)$.

\paragraph{What the defender has done.}
The defender has given each agent a role, and we assume those roles are scoped
by the usual advice: agent~A's policies do not name agent~B's data. Where A
can read B's data directly there is no escalation to find, and an existing
single-identity tool reports it.

\paragraph{Out of scope.}
The command route of Section~\ref{sec:background} remains available to the
attacker throughout. We do not model the host, the container runtime, the platform's isolation
between tenants, or a compromise of the cloud control plane. An agent reading
another's output through a channel the deployment intended is not an
escalation here, because Equation~\ref{eq:esc} puts that reach inside
$\auth$.

\paragraph{Why a single-identity analysis misses this.}
The escalation is a property of the pair, and it appears only when the analysis
follows the edge from A to B and then continues with B's permissions. That edge
is the subject of Section~\ref{sec:approach}.
